\section{What a human should take from this}
\label{sec:playbook}

FishBot v0.4 maintains a $54 \times 6$ posterior and evaluates a hundred-odd
candidate asks per turn. A person cannot. But the mechanisms the ablations
identify as decisive are not the expensive ones, and three of them are cheap
enough to run at a table.

\subsection{The deduction that pays for itself}

\textbf{An ask is a confession.} When somebody asks for the queen of hearts, they
have told the table two things: they do not hold the queen, and they \emph{do}
hold at least one other high heart. The second is the one people forget, and
together with card counting it is where essentially all of the agent's strength
lives: strip either from FishBot and it drops from winning roughly three games in
four to losing three in four. It applies to teammates as
much as to opponents, and it does not expire: the cards it refers to cannot move
except by a public ask.

\textbf{A miss is a second confession.} It rules the named card out of the target
permanently, for the same reason.

Keep these as exclusions, not impressions. ``Somebody wanted low clubs'' is worth
little; ``the 7 of clubs is not with Iris, not with Theo, and Sage holds one of
the other five'' is worth a great deal, because it combines with counting.

\subsection{Counting closes the deductions}

Card counts are public, and they are the constraint that turns partial exclusions
into certainties. If a player has four cards and you have located four of them,
they hold nothing else --- every unresolved card is excluded from them at once.
Run this after every transfer, on the two players whose counts changed. Most of
FishBot's ``impossible'' deductions come from nothing more sophisticated than
this.

\subsection{The claim you should not make}

\textbf{If your team holds all six cards of a half-suit, nobody can take it.} An
opponent needs a card of a half-suit to ask in it, and they have none. The set is
frozen until it is claimed (Theorem~\ref{thm:locked}), and if an opponent claims
it they must be wrong, which hands it to you. It follows that claiming early
buys you no safety --- and costs you a little, because the six cards leaving play
sharpens everyone's count of what remains. So:

Two qualifications, because the simulator does not simply endorse patience. The
fitted agent is in fact the \emph{less} patient of the two policies we measured
(\S\ref{sec:theory}), and the ``cash locked half-suits immediately'' switch is
inert in the shipped configuration, so no experiment here isolates the value of
waiting. What the rules do guarantee is the safety, not the profit.

\begin{itemize}
  \item If you are sure your side has all six but unsure which teammate has what,
        you are under no time pressure --- nobody can take it. Use that freedom
        to resolve the split rather than guessing, but do not sit on it out of
        principle: the cards buy you nothing at the table, and the point only
        scores once you claim.
  \item If a half-suit is \emph{not} yet fully yours, the calculus reverses:
        waiting risks an opponent asking one of your cards away, so a confident
        claim is right.
  \item Do not sit on a hand that is entirely locked cards. Those cards license
        only asks that are certain to fail, so your turn is worth nothing.
        Cash the half-suit, empty your hand, and hand the turn to a teammate who
        can use it.
\end{itemize}

\subsection{Choosing the ask}

Rank candidates roughly in this order, which is what the fitted weights do:

\begin{enumerate}
  \item \textbf{Certainty first.} A card you have located is free: you take it and
        you keep the turn. Take every one of these before anything speculative.
  \item \textbf{Then the run, not the card.} Because a hit keeps the turn, the
        value of an ask is the length of the sequence it starts. Prefer a likely
        card that leaves another likely card behind it over an equally likely
        card that leaves nothing.
  \item \textbf{Then the lock.} An ask that would complete your team's ownership
        of a half-suit is worth more than its hit probability suggests, because
        it converts a contested set into a frozen one.
  \item \textbf{Then, who receives the miss.} You are choosing your successor as
        much as your card. Avoid handing the turn to the player who is visibly
        loaded in a half-suit where your side is exposed.
  \item \textbf{And what you are giving away.} Your first ask in a half-suit
        announces that you hold one of its cards. If you are choosing between two
        near-equal asks, prefer the one in a half-suit you have already revealed.
\end{enumerate}

\subsection{Declaring, concretely}

Before you claim, say the six cards and a name for each, out loud in your head.
Separate ``our side has it'' from ``I know which of us has it'' --- they are
different claims and only the second is enough. If either is shaky and the
half-suit is already locked, wait; if it is not locked, weigh the claim against
the chance an opponent takes a card first.

\subsection{The endgame nobody plans for}

When one team runs out of cards, the other must claim every remaining half-suit
with no more asking. Two consequences follow that most players get wrong.

First, \textbf{running out of cards is not a defeat.} A cardless player cannot be
asked, so nothing can be taken from them, and a cardless \emph{team} forces the
opponents to declare everything --- which they will sometimes get wrong, handing
you the set. Both of these are consequences of the rules rather than findings of
this study: we did not isolate the value of emptying a hand deliberately, and we
flag it as an option worth knowing rather than a measured recommendation.

Second, when it is your team that must claim everything, \textbf{claim your
surest half-suit first.} A correct claim announces the exact split of six cards,
which changes everyone's counts and can settle the half-suit you were unsure of.
Order matters, and the safe claim first is the right order.

\subsection{What the simulator does not support}

The v0.3 study found no reliable benefit from deliberate misdirection, and
nothing in this study changes that. FishBot v0.4 has no bluffing rule and beats
the misdirection archetype essentially always. Theatre that does not change what
you know or who holds the turn is not worth paying for.
